\section{PL 条件让部分非凸问题恢复全局速率}
\label{sec:09-Convex-Optimization/08-Nonconvex-and-Stochastic-Optimization/03}

\subsection{没有凸性，是不是就彻底没办法了？}

凸优化的全局收敛率有一个令人安心的逻辑链：凸性 \(\to\) 切平面在函数图像下方 \(\to\) 梯度大小与函数值差距直接挂钩 \(\to\) 每次梯度步按固定比例缩小 gap \(\to\) 线性收敛。每步都走得有根有据。

但很多真实问题不凸。深度学习的目标函数充满鞍点、弯曲和谷地。凸性这个前提一旦撤掉，上述逻辑链的第一环就断了——切平面下界不再成立。那么，线性收敛这种强结论是不是也一起被判了死刑？

这是整个非凸优化理论要回答的核心问题。本节介绍的 PL 条件给出了一个出乎意料的答案：部分非凸问题可以恢复凸性级别的收敛保证，路径不是修补切平面下界，而是绕过去，直接建立梯度和函数值之间的不等式。

\subsection{PL 条件：一句简单到让人怀疑的不等式}

Polyak--\L{}ojasiewicz（PL）条件是这么写的：存在常数 \(\mu>0\)，使得对所有 \(x\) 都有

\[
\frac12 \norm{\nabla f(x)}^2 \ge \mu\bigl(f(x) - f^\star\bigr),
\]

其中 \(f^\star = \inf_x f(x)\) 是全局最优值。

它只在说一件事：只要当前位置的函数值还明显高于最优值，梯度就不能任意小。你想停在一个梯度接近零但函数值很差的地方？PL 条件不允许这种情况出现——因为梯度小意味着 \(f(x)\) 已经靠近 \(f^\star\)。

这个条件的威力在于它直接排除了非凸优化中两个最常见的失败模式：一是``梯度很小但你离最优还很远''的平坦高原；二是上节讨论的``坏驻点''——它们梯度为零却根本不是极小。PL 条件下，任何梯度为零的点自动满足 \(f(x)=f^\star\)，函数值高于 \(f^\star\) 的驻点根本不存在。

\subsection{把 PL 接到下降引理上，线性收敛自己就掉出来了}

如果 \(f\) 同时满足 \(L\)-光滑（梯度 Lipschitz 连续），取步长为 \(1/L\) 的梯度下降步的标准下降引理给出：

\[
f(x_{k+1}) \le f(x_k) - \frac{1}{2L}\norm{\nabla f(x_k)}^2.
\]

这个不等式不依赖凸性，只依赖光滑性。现在把 PL 条件的不等式代入右边：

\[
f(x_{k+1}) - f^\star
\le \bigl(f(x_k) - f^\star\bigr) - \frac{\mu}{L}\bigl(f(x_k) - f^\star\bigr)
= \left(1 - \frac{\mu}{L}\right)\bigl(f(x_k) - f^\star\bigr).
\]

递推 \(k\) 步：

\[
f(x_k) - f^\star \le \left(1 - \frac{\mu}{L}\right)^k \bigl(f(x_0) - f^\star\bigr).
\]

这和强凸函数的线性收敛形式完全相同，却从未用到凸性的弦不等式或 Hessian 处处半正定的要求。仅凭光滑性和 PL 这两条，就把梯度下降的全局线性收敛率建立了起来。

\subsection{PL 比强凸弱，而且根本不要求凸}

强凸函数必定满足 PL 条件（强凸直接给出了 \(\norm{\nabla f(x)}^2 \ge 2\mu(f(x)-f^\star)\) 的形式）。反向不成立：满足 PL 的函数可以不强凸，甚至可以不满秩 Hessian。

考虑秩亏最小二乘：

\[
f(x) = \frac12 \norm{Ax - b}^2.
\]

如果 \(A\) 的列不满秩，最优解构成一个仿射子空间，而不是唯一的点——这样的函数不可能是强凸的（强凸要求唯一极小）。但 \(A\) 的非零奇异值仍然提供 PL 型不等式，因为梯度 \(\nabla f(x)=A^{\trans}(Ax-b)\) 的范数与 \(Ax\) 到 \(b\) 的距离直接相关，而该距离又控制着函数值到最优值的 gap。强凸性缺位，PL 却安然在岗。

更令人意外的是，PL 甚至可以覆盖非凸函数。构造一个具体例子：

\[
f(x) = x^2 + 3(1-\cos x).
\]

它的二阶导数

\[
f''(x) = 2 + 3\cos x
\]

在 \(x\) 靠近 \(\pi\) 时为负——函数在该区域非凸。但 \(f\) 只有一个驻点 \(x=0\)（全局极小），而且可以验证以下两个不等式：

\[
|f'(x)| \ge |x|,\qquad
f(x) \le \frac52 x^2.
\]

第一条的证明分区间：在 \(0\le x\le\pi\) 利用 \(\sin x\ge 0\)；在 \(x\ge\pi\) 利用 \(2x+3\sin x\ge 2x-3\ge x\)；负半轴由对称性可得。第二条来自 \(1-\cos x\le x^2/2\)。

两条联立：

\[
\frac12 |f'(x)|^2 \ge \frac12 x^2 \ge \frac15 f(x),
\]

即 \(f\) 满足 \(\mu=1/5\) 的 PL 条件。一个货真价实的非凸函数，同样享有梯度下降的全局线性收敛率。

\subsection{PL 干掉了坏驻点，但剩下的事情不简单}

在 PL 条件下，\(\nabla f(x)=0\) 立刻推出 \(f(x)=f^\star\)，所有驻点都是全局最优。这一点确实漂亮。但它不保证最优解唯一——可能有一整片点都到达 \(f^\star\)，算法收敛到哪一点取决于初始化和路径。PL 也不给出参数空间中的距离收缩，只给出函数值的 gap 缩减。

更大的警示在这里：PL 是目标函数的全局结构假设，不是可以从训练曲线上反推的经验事实。你看到损失下降了，不能倒过来说``所以这个函数满足 PL''。一般深度网络在整个参数空间上是否具有统一的 PL 常数，没有理论保证。某些过参数化模型可能在特定初始化邻域、特定数据分布或沿训练轨迹的局部区域出现类似 PL 的行为，但这个结论仅在限定的范围内成立，不能无条件推广。

PL 的真正价值不在于宣称``深度学习满足 PL''，而在于展示了一条中间道路的存在：问题不凸，不代表必须彻底放弃全局收敛理论。但要恢复强结论，你必须找到某种可验证的替代结构——PL 是这类替代结构中门槛最低、后果最直接的一种。它的存在提醒我们：凸性和线性收敛之间的那道墙，并不如第一眼看上去那么牢固。
